
\section*{Part III. Atomic physics \hfill 12 / 15}
\subsection*{1. Fine structure and hyperfine structure}
In a spin--orbit model, \(H_{\rm SO}\propto\mathbf L\cdot\mathbf S\).
Since \(\mathbf J=\mathbf L+\mathbf S\),
\[
\langle\mathbf L\cdot\mathbf S\rangle
=\frac{\hbar^2}{2}
[J(J+1)-L(L+1)-S(S+1)].
\]
This separates \(J\) levels within a fixed LS term.
For the magnetic-dipole hyperfine interaction,
\[
H_{\rm hfs}=\frac A{\hbar^2}\mathbf I\cdot\mathbf J,\qquad
\boxed{E_F=\frac A2[F(F+1)-I(I+1)-J(J+1)].}
\]
Here \(F=|I-J|,\ldots,I+J\); each level has
\(2F+1\) projections when no magnetic field breaks degeneracy.
Electric-quadrupole hyperfine terms can also appear for
sufficient nuclear/electronic angular momentum.

\textbf{Worked example:} \(I=5/2,J=1/2\) gives \(F=2,3\).
Using \(I(I+1)=35/4,\ J(J+1)=3/4\),
\[
\boxed{E_{2}=-\frac{7A}{4},\quad
E_{3}=\frac{5A}{4},\quad \Delta E=3A.}
\]
The multiplicities \(5+7=12=(2I+1)(2J+1)\).
For \(A<0\), the energy order is reversed.

\subsection*{2. Magnetic and electric fields}
In a weak magnetic field, \(F,m_F\) are approximately good
quantum numbers. The linear Zeeman shift is
\[
\Delta E\simeq g_F\mu_Bm_FB,\quad
g_F\simeq g_J
\frac{F(F+1)+J(J+1)-I(I+1)}{2F(F+1)}
\]
for \(F>0\) and neglecting the small nuclear magnetic term.
At strong fields, \(I\) and \(J\) decouple approximately,
so a basis \(\ket{Jm_J,Im_I}\) becomes useful.
A static electric field induces a Stark shift. For a
nondegenerate parity eigenstate with no permanent dipole,
\[
\Delta E_{\rm Stark}^{(2)}
=-\tfrac12\alpha(0)\mathcal E^2,\qquad
\alpha(0)=2\sum_{n\ne i}
\frac{|\langle n|d_z|i\rangle|^2}{E_n-E_i}
\]
with field along \(z\) (and compatible dipole conventions).
Degenerate states require diagonalisation rather than this
simple nondegenerate formula.

\subsection*{3. Isotope shifts}
Changing neutron number changes the nuclear mass and charge
distribution. Spectral isotope shifts reflect normal/specific
mass shifts and field (volume) shifts; these can reveal nuclear
radius differences when atomic sensitivity coefficients are known.
\begin{methodbox}
\textbf{Method:} choose the coupling/field regime,
list allowed \(F\), compute splitting; for a perturbation,
identify the operator and apply first- or second-order theory.
\end{methodbox}
\src{Week 8 slides on relativistic structure; Week 9 pp.~2--7, 21--27; hyperfine tutorial and transcripts.}
